% Appendix A: Plagiarism Report
\chapter{Plagiarism Report}
\section{Plagiarism Check Results}
The plagiarism check for this project report was conducted using an approved anti-plagiarism tool prior to submission. The report confirms that the content of this work is original and within the acceptable similarity threshold stipulated by RV College of Engineering.

\begin{figure}[H]
    \centering
    % Placeholder for plagiarism report screenshot/scan
    % Replace with: \includegraphics[width=0.9\textwidth]{Figures/plagiarism_report}
    \fbox{\parbox{0.85\textwidth}{\centering \vspace{3cm} \large\textit{[Plagiarism Report Screenshot / Scan to be inserted here]} \vspace{3cm}}}
    \makeatletter
    \let\saved@addcontentsline\addcontentsline
    \renewcommand{\addcontentsline}[3]{}
    \caption{Plagiarism Check Report}
    \let\addcontentsline\saved@addcontentsline
    \makeatother
    \label{fig:plagiarism_report}
\end{figure}

\textbf{Note:} The overall similarity index is within the permissible limit. All citations and references have been properly acknowledged in the Bibliography section of this report.

\clearpage

% Appendix B: Source Code
\chapter{Source Code}
\section{Key Module Implementations}
This appendix presents representative code excerpts from the key modules of the Agentic Supply-Chain Digital Twin system. The full source code is available in the project's GitHub repository.

\subsection{Friction Sentry Module}
The following listing presents the core shock parsing logic of the Friction Sentry module, which maps raw disruption signals to structured risk parameters.

\begin{tcblisting}{listing only,colback=gray!10!white, breakable, boxsep=0pt,top=1mm,bottom=1mm,left=1mm,right=1mm,
listing options={
language=Python,
basicstyle=\small\ttfamily,
keywordstyle=\color{blue},
commentstyle=\color{gray},
stringstyle=\color{red!70!black},
backgroundcolor=\color{gray!10},
showstringspaces=false,
numbers=left,
numberstyle=\tiny\color{gray},
numbersep=5pt}}
# friction_sentry.py - Core Signal Parsing Logic
import json

def parse_shock_signal(location: str, signal_type: str,
                        severity: str, rules_db: dict) -> dict:
    """
    Maps an unstructured disruption signal to a structured
    shock payload using the semantic mapping rules library.
    """
    # Resolve location to system node ID
    node_id = rules_db["location_mapping"].get(
        location.lower().replace(" ", "_"),
        "unknown_node"
    )

    # Extract signal-specific multipliers
    signal_rules = rules_db["signal_type_rules"].get(
        signal_type, {}
    )
    severity_data = signal_rules.get(severity, {})

    # Build structured shock payload
    shock_payload = {
        "target_node": node_id,
        "signal_type": signal_type,
        "severity_multiplier": severity_data.get(
            "cost_multiplier", 1.0
        ),
        "delay_days": severity_data.get("delay_days", 0),
        "cascade_nodes": severity_data.get("cascade_nodes", [])
    }
    return shock_payload
\end{tcblisting}

\subsection{Dijkstra Routing Engine}
The following listing demonstrates the core pathfinding logic used to compute optimal alternate routes under shocked network conditions.

\begin{tcblisting}{listing only,colback=gray!10!white, breakable, boxsep=0pt,top=1mm,bottom=1mm,left=1mm,right=1mm,
listing options={
language=Python,
basicstyle=\small\ttfamily,
keywordstyle=\color{blue},
commentstyle=\color{gray},
stringstyle=\color{red!70!black},
backgroundcolor=\color{gray!10},
showstringspaces=false,
numbers=left,
numberstyle=\tiny\color{gray},
numbersep=5pt}}
# engine/dijkstra_router.py - Optimal Path Computation
import networkx as nx

def dijkstra_best_path(graph: nx.DiGraph,
                        source: str,
                        target: str,
                        optimize_by: str = "cost") -> dict:
    """
    Computes the shortest path from source to target
    on the dynamically weighted logistics graph.
    """
    weight_attr = (
        "freight_cost" if optimize_by == "cost"
        else "transit_time"
    )

    try:
        path = nx.dijkstra_path(
            graph, source, target, weight=weight_attr
        )
        total_cost = sum(
            graph[u][v].get("freight_cost", 0)
            for u, v in zip(path, path[1:])
        )
        total_days = sum(
            graph[u][v].get("transit_time", 0)
            for u, v in zip(path, path[1:])
        )
        return {
            "path": path,
            "total_cost_usd": round(total_cost, 2),
            "total_transit_days": round(total_days, 2),
            "status": "optimal_route_found"
        }
    except nx.NetworkXNoPath:
        return {"status": "no_path_available", "path": []}
\end{tcblisting}

\subsection{CrewAI Agent Orchestration}
The following listing presents the multi-agent orchestration setup using the CrewAI framework, showing agent definitions and task chaining.

\begin{tcblisting}{listing only,colback=gray!10!white, breakable, boxsep=0pt,top=1mm,bottom=1mm,left=1mm,right=1mm,
listing options={
language=Python,
basicstyle=\small\ttfamily,
keywordstyle=\color{blue},
commentstyle=\color{gray},
stringstyle=\color{red!70!black},
backgroundcolor=\color{gray!10},
showstringspaces=false,
numbers=left,
numberstyle=\tiny\color{gray},
numbersep=5pt}}
# agent_orchestrator.py - CrewAI Multi-Agent Setup
from crewai import Agent, Task, Crew, Process

def run_agentic_pipeline(shock_context: dict,
                          llm) -> dict:
    """Orchestrates the 3-agent CrewAI pipeline."""

    # Agent 1: Logistics Optimizer
    optimizer = Agent(
        role="Logistics Optimizer",
        goal="Find the optimal alternate shipping route.",
        backstory="An expert in maritime routing algorithms.",
        llm=llm,
        tools=[DijkstraOptimalRouteTool()]
    )

    # Agent 2: Cost & Strategy Analyst
    analyst = Agent(
        role="Cost and Strategy Analyst",
        goal="Verify inventory health and evaluate costs.",
        backstory="A supply chain economist.",
        llm=llm,
        tools=[CheckInventoryTool(), EmergencyRestockTool()]
    )

    # Agent 3: SME Communicator
    communicator = Agent(
        role="SME Communicator",
        goal="Compile findings into an executive summary.",
        backstory="A logistics communications expert.",
        llm=llm
    )

    # Task Chaining with context passing
    t1 = Task(description=f"Find optimal route given: "
              f"{shock_context}", agent=optimizer)
    t2 = Task(description="Evaluate inventory and costs.",
              agent=analyst, context=[t1])
    t3 = Task(description="Generate executive summary.",
              agent=communicator, context=[t1, t2])

    crew = Crew(
        agents=[optimizer, analyst, communicator],
        tasks=[t1, t2, t3],
        process=Process.sequential
    )
    return crew.kickoff()
\end{tcblisting}